\textbf{Những công việc đã thực hiện trong khóa luận}

Khóa luận đã nghiên cứu và phát triển một nguyên mẫu hệ thống kiểm soát chất lượng không khí trong căn hộ chung cư có tích hợp cây lan ý và công nghệ IoT. Trên cơ sở tham khảo các hướng tiếp cận về IAQ, cây trồng/thủy canh và hệ thống giám sát thông minh, đề tài xây dựng mô hình phù hợp với bối cảnh căn hộ: quy mô nhỏ, chi phí hợp lý, dễ lắp đặt, có khả năng theo dõi dữ liệu theo thời gian thực và hỗ trợ điều khiển thiết bị. Hệ thống được tổ chức theo kiến trúc IoT gồm node cảm biến Arduino Pro Mini, liên kết LoRa, gateway ESP32, nền tảng ThingsBoard và ứng dụng Flutter.

Các kết quả chính đã đạt được gồm:
\begin{itemize}
    \item Thiết kế node cảm biến đọc được các thông số TDS, pH, nhiệt độ, độ ẩm, CO2 và ánh sáng.
    \item Xây dựng cơ chế truyền dữ liệu không dây giữa node và gateway bằng hai module LoRa.
    \item Phát triển gateway ESP32 để nhận dữ liệu LoRa, xử lý gói tin và gửi dữ liệu lên ThingsBoard.
    \item Thiết kế mô hình dữ liệu trên ThingsBoard phục vụ hiển thị, cảnh báo và điều khiển.
    \item Xây dựng ứng dụng Flutter để người dùng theo dõi thông số môi trường và điều khiển thiết bị.
    \item Triển khai điều khiển relay bật/tắt quạt, relay bật/tắt đèn và điều khiển điều hòa bằng IR tại node.
\end{itemize}

Kết quả triển khai cho thấy hệ thống có thể giám sát môi trường trong nhà, hỗ trợ chăm sóc cây lan ý và thực hiện điều khiển thiết bị từ xa. Kiến trúc node - gateway - cloud - app phù hợp với điều kiện căn hộ chung cư và có thể mở rộng khi bổ sung thêm node cảm biến. Tuy nhiên, để đánh giá hiệu quả cải thiện IAQ ở mức định lượng, hệ thống cần được kiểm thử trong thời gian dài hơn, với cảm biến được hiệu chuẩn, nhiều điều kiện thông gió khác nhau và các thông số bổ sung như VOC hoặc bụi mịn.
